\begin{abstract}
El melanoma se parece a un lunar común en sus comienzos, y la regla dermatoscópica ABCD ayuda a distinguirlos a partir de la asimetría, el borde, los colores y las estructuras de la lesión. En este trabajo se desarrolló un algoritmo que mide esos criterios en las 200 imágenes de la base PH2 con procesamiento clásico de imágenes, sin entrenar un clasificador. Las imágenes se separaron en 140 de entrenamiento y 60 de prueba, y todo parámetro se fijó con las primeras. La segmentación por crecimiento de regiones alcanzó un Dice medio de 0,916 en prueba. Quitar los pelos con un cierre morfológico no mejoró la segmentación, por lo que la primera hipótesis no se confirmó. El criterio D, estimado detectando cada estructura por separado, separó en prueba las lesiones sospechosas de los lunares comunes con un AUC de 0,67 (IC 95\,\%, 0,54 a 0,80), y la diferencia de saturación entre la periferia y el centro de la lesión, con 0,71 (0,58 a 0,84). La asimetría y los colores medidos coincidieron con los del dermatólogo con un $\kappa$ de 0,38 y 0,24.
El TDS automático separó las lesiones sospechosas de los nevos comunes con un AUC de 0,77 (0,65 a 0,89) y, con un umbral elegido en entrenamiento, una sensibilidad de 0,69 y una especificidad de 0,88. Agregarle términos de zonas no mejoró esa separación de forma significativa, por lo que la segunda hipótesis se confirmó sólo en parte.
\end{abstract}
